\IfFileExists{stacks-project.cls}{%
\documentclass{stacks-project}
}{%
\documentclass{amsart}
}

% For dealing with references we use the comment environment
\usepackage{verbatim}
\newenvironment{reference}{\comment}{\endcomment}
%\newenvironment{reference}{}{}
\newenvironment{slogan}{\comment}{\endcomment}
\newenvironment{history}{\comment}{\endcomment}

% For commutative diagrams we use Xy-pic
\usepackage[all]{xy}

% We use 2cell for 2-commutative diagrams.
\xyoption{2cell}
\UseAllTwocells

% We use multicol for the list of chapters between chapters
\usepackage{multicol}

% This is generally recommended for better output
\usepackage{lmodern}
\usepackage[T1]{fontenc}

% For cross-file-references
\usepackage{xr-hyper}

% Package for hypertext links:
\usepackage{hyperref}

% For any local file, say "hello.tex" you want to link to please
% use \externaldocument[hello-]{hello}
\externaldocument[introduction-]{introduction}
\externaldocument[conventions-]{conventions}
\externaldocument[sets-]{sets}
\externaldocument[categories-]{categories}
\externaldocument[topology-]{topology}
\externaldocument[sheaves-]{sheaves}
\externaldocument[sites-]{sites}
\externaldocument[stacks-]{stacks}
\externaldocument[fields-]{fields}
\externaldocument[algebra-]{algebra}
\externaldocument[brauer-]{brauer}
\externaldocument[homology-]{homology}
\externaldocument[derived-]{derived}
\externaldocument[simplicial-]{simplicial}
\externaldocument[more-algebra-]{more-algebra}
\externaldocument[smoothing-]{smoothing}
\externaldocument[modules-]{modules}
\externaldocument[sites-modules-]{sites-modules}
\externaldocument[injectives-]{injectives}
\externaldocument[cohomology-]{cohomology}
\externaldocument[sites-cohomology-]{sites-cohomology}
\externaldocument[dga-]{dga}
\externaldocument[dpa-]{dpa}
\externaldocument[sdga-]{sdga}
\externaldocument[hypercovering-]{hypercovering}
\externaldocument[schemes-]{schemes}
\externaldocument[constructions-]{constructions}
\externaldocument[properties-]{properties}
\externaldocument[morphisms-]{morphisms}
\externaldocument[coherent-]{coherent}
\externaldocument[divisors-]{divisors}
\externaldocument[limits-]{limits}
\externaldocument[varieties-]{varieties}
\externaldocument[topologies-]{topologies}
\externaldocument[descent-]{descent}
\externaldocument[perfect-]{perfect}
\externaldocument[more-morphisms-]{more-morphisms}
\externaldocument[flat-]{flat}
\externaldocument[groupoids-]{groupoids}
\externaldocument[more-groupoids-]{more-groupoids}
\externaldocument[etale-]{etale}
\externaldocument[chow-]{chow}
\externaldocument[intersection-]{intersection}
\externaldocument[pic-]{pic}
\externaldocument[weil-]{weil}
\externaldocument[adequate-]{adequate}
\externaldocument[dualizing-]{dualizing}
\externaldocument[duality-]{duality}
\externaldocument[discriminant-]{discriminant}
\externaldocument[derham-]{derham}
\externaldocument[local-cohomology-]{local-cohomology}
\externaldocument[algebraization-]{algebraization}
\externaldocument[curves-]{curves}
\externaldocument[resolve-]{resolve}
\externaldocument[models-]{models}
\externaldocument[functors-]{functors}
\externaldocument[equiv-]{equiv}
\externaldocument[pione-]{pione}
\externaldocument[etale-cohomology-]{etale-cohomology}
\externaldocument[proetale-]{proetale}
\externaldocument[relative-cycles-]{relative-cycles}
\externaldocument[more-etale-]{more-etale}
\externaldocument[trace-]{trace}
\externaldocument[crystalline-]{crystalline}
\externaldocument[spaces-]{spaces}
\externaldocument[spaces-properties-]{spaces-properties}
\externaldocument[spaces-morphisms-]{spaces-morphisms}
\externaldocument[decent-spaces-]{decent-spaces}
\externaldocument[spaces-cohomology-]{spaces-cohomology}
\externaldocument[spaces-limits-]{spaces-limits}
\externaldocument[spaces-divisors-]{spaces-divisors}
\externaldocument[spaces-over-fields-]{spaces-over-fields}
\externaldocument[spaces-topologies-]{spaces-topologies}
\externaldocument[spaces-descent-]{spaces-descent}
\externaldocument[spaces-perfect-]{spaces-perfect}
\externaldocument[spaces-more-morphisms-]{spaces-more-morphisms}
\externaldocument[spaces-flat-]{spaces-flat}
\externaldocument[spaces-groupoids-]{spaces-groupoids}
\externaldocument[spaces-more-groupoids-]{spaces-more-groupoids}
\externaldocument[bootstrap-]{bootstrap}
\externaldocument[spaces-pushouts-]{spaces-pushouts}
\externaldocument[spaces-chow-]{spaces-chow}
\externaldocument[groupoids-quotients-]{groupoids-quotients}
\externaldocument[spaces-more-cohomology-]{spaces-more-cohomology}
\externaldocument[spaces-simplicial-]{spaces-simplicial}
\externaldocument[spaces-duality-]{spaces-duality}
\externaldocument[formal-spaces-]{formal-spaces}
\externaldocument[restricted-]{restricted}
\externaldocument[spaces-resolve-]{spaces-resolve}
\externaldocument[formal-defos-]{formal-defos}
\externaldocument[defos-]{defos}
\externaldocument[cotangent-]{cotangent}
\externaldocument[examples-defos-]{examples-defos}
\externaldocument[algebraic-]{algebraic}
\externaldocument[examples-stacks-]{examples-stacks}
\externaldocument[stacks-sheaves-]{stacks-sheaves}
\externaldocument[criteria-]{criteria}
\externaldocument[artin-]{artin}
\externaldocument[quot-]{quot}
\externaldocument[stacks-properties-]{stacks-properties}
\externaldocument[stacks-morphisms-]{stacks-morphisms}
\externaldocument[stacks-limits-]{stacks-limits}
\externaldocument[stacks-cohomology-]{stacks-cohomology}
\externaldocument[stacks-perfect-]{stacks-perfect}
\externaldocument[stacks-introduction-]{stacks-introduction}
\externaldocument[stacks-more-morphisms-]{stacks-more-morphisms}
\externaldocument[stacks-geometry-]{stacks-geometry}
\externaldocument[moduli-]{moduli}
\externaldocument[moduli-curves-]{moduli-curves}
\externaldocument[examples-]{examples}
\externaldocument[exercises-]{exercises}
\externaldocument[guide-]{guide}
\externaldocument[desirables-]{desirables}
\externaldocument[coding-]{coding}
\externaldocument[obsolete-]{obsolete}
\externaldocument[fdl-]{fdl}
\externaldocument[index-]{index}

% Theorem environments.
%
\theoremstyle{plain}
\newtheorem{theorem}[subsection]{Theorem}
\newtheorem{proposition}[subsection]{Proposition}
\newtheorem{lemma}[subsection]{Lemma}

\theoremstyle{definition}
\newtheorem{definition}[subsection]{Definition}
\newtheorem{example}[subsection]{Example}
\newtheorem{exercise}[subsection]{Exercise}
\newtheorem{situation}[subsection]{Situation}

\theoremstyle{remark}
\newtheorem{remark}[subsection]{Remark}
\newtheorem{remarks}[subsection]{Remarks}

\numberwithin{equation}{subsection}

% Macros
%
\def\lim{\mathop{\mathrm{lim}}\nolimits}
\def\colim{\mathop{\mathrm{colim}}\nolimits}
\def\Spec{\mathop{\mathrm{Spec}}}
\def\Hom{\mathop{\mathrm{Hom}}\nolimits}
\def\Ext{\mathop{\mathrm{Ext}}\nolimits}
\def\SheafHom{\mathop{\mathcal{H}\!\mathit{om}}\nolimits}
\def\SheafExt{\mathop{\mathcal{E}\!\mathit{xt}}\nolimits}
\def\Sch{\mathit{Sch}}
\def\Mor{\mathop{\mathrm{Mor}}\nolimits}
\def\Ob{\mathop{\mathrm{Ob}}\nolimits}
\def\Sh{\mathop{\mathit{Sh}}\nolimits}
\def\NL{\mathop{N\!L}\nolimits}
\def\CH{\mathop{\mathrm{CH}}\nolimits}
\def\proetale{{pro\text{-}\acute{e}tale}}
\def\etale{{\acute{e}tale}}
\def\QCoh{\mathit{QCoh}}
\def\Ker{\mathop{\mathrm{Ker}}}
\def\Im{\mathop{\mathrm{Im}}}
\def\Coker{\mathop{\mathrm{Coker}}}
\def\Coim{\mathop{\mathrm{Coim}}}

% Boxtimes
%
\DeclareMathSymbol{\boxtimes}{\mathbin}{AMSa}{"02}

%
% Macros for moduli stacks/spaces
%
\def\QCohstack{\mathcal{QC}\!\mathit{oh}}
\def\Cohstack{\mathcal{C}\!\mathit{oh}}
\def\Spacesstack{\mathcal{S}\!\mathit{paces}}
\def\Quotfunctor{\mathrm{Quot}}
\def\Hilbfunctor{\mathrm{Hilb}}
\def\Curvesstack{\mathcal{C}\!\mathit{urves}}
\def\Polarizedstack{\mathcal{P}\!\mathit{olarized}}
\def\Complexesstack{\mathcal{C}\!\mathit{omplexes}}
% \Pic is the operator that assigns to X its picard group, usage \Pic(X)
% \Picardstack_{X/B} denotes the Picard stack of X over B
% \Picardfunctor_{X/B} denotes the Picard functor of X over B
\def\Pic{\mathop{\mathrm{Pic}}\nolimits}
\def\Picardstack{\mathcal{P}\!\mathit{ic}}
\def\Picardfunctor{\mathrm{Pic}}
\def\Deformationcategory{\mathcal{D}\!\mathit{ef}}
\usepackage{cleveref}
\crefname{originalreference}{original reference}{original references}
\Crefname{originalreference}{Original reference}{Original references}
\makeatletter
\newlabel{algebra-lemma-finite-gl-dim-finite-dim-regular}{{00OE}{source}{}{sourcepointer-1}{}}
\newlabel{algebra-lemma-finite-gl-dim-finite-dim-regular@cref}{{[originalreference][0][]00OE}{[1][0][]source}{}{}{}}
\makeatother
\crefname{originalreference}{original reference}{original references}
\Crefname{originalreference}{Original reference}{Original references}
\makeatletter
\newlabel{cohomology-lemma-pseudo-coherent-pullback}{{09U7}{source}{}{sourcepointer-2}{}}
\newlabel{cohomology-lemma-pseudo-coherent-pullback@cref}{{[originalreference][0][]09U7}{[1][0][]source}{}{}{}}
\makeatother
\crefname{originalreference}{original reference}{original references}
\Crefname{originalreference}{Original reference}{Original references}
\makeatletter
\newlabel{cohomology-lemma-tensor-pseudo-coherent}{{09J3}{source}{}{sourcepointer-3}{}}
\newlabel{cohomology-lemma-tensor-pseudo-coherent@cref}{{[originalreference][0][]09J3}{[1][0][]source}{}{}{}}
\makeatother
\crefname{originalreference}{original reference}{original references}
\Crefname{originalreference}{Original reference}{Original references}
\makeatletter
\newlabel{cohomology-lemma-tor-amplitude-stalk}{{09U9}{source}{}{sourcepointer-4}{}}
\newlabel{cohomology-lemma-tor-amplitude-stalk@cref}{{[originalreference][0][]09U9}{[1][0][]source}{}{}{}}
\makeatother
\crefname{originalreference}{original reference}{original references}
\Crefname{originalreference}{Original reference}{Original references}
\makeatletter
\newlabel{cohomology-lemma-variant-derived-pullback}{{08DE}{source}{}{sourcepointer-5}{}}
\newlabel{cohomology-lemma-variant-derived-pullback@cref}{{[originalreference][0][]08DE}{[1][0][]source}{}{}{}}
\makeatother
\crefname{originalreference}{original reference}{original references}
\Crefname{originalreference}{Original reference}{Original references}
\makeatletter
\newlabel{more-algebra-lemma-complex-finite-tor-dimension-modules}{{066H}{source}{}{sourcepointer-6}{}}
\newlabel{more-algebra-lemma-complex-finite-tor-dimension-modules@cref}{{[originalreference][0][]066H}{[1][0][]source}{}{}{}}
\makeatother
\crefname{originalreference}{original reference}{original references}
\Crefname{originalreference}{Original reference}{Original references}
\makeatletter
\newlabel{perfect-lemma-direct-image-coherent}{{08E2}{source}{}{sourcepointer-7}{}}
\newlabel{perfect-lemma-direct-image-coherent@cref}{{[originalreference][0][]08E2}{[1][0][]source}{}{}{}}
\makeatother
\crefname{originalreference}{original reference}{original references}
\Crefname{originalreference}{Original reference}{Original references}
\makeatletter
\newlabel{perfect-lemma-identify-pseudo-coherent-noetherian}{{08E8}{source}{}{sourcepointer-8}{}}
\newlabel{perfect-lemma-identify-pseudo-coherent-noetherian@cref}{{[originalreference][0][]08E8}{[1][0][]source}{}{}{}}
\makeatother
\title{Fourier-Mukai functor preserves bounded coherent in four cases}
\author{The Stacks Project Authors}
\date{stacks-project:equiv.tex:local-snapshot-20261012; GNU FDL 1.2 or later}
\begin{document}
\maketitle
\noindent\textit{Allowed standard prerequisites (editorial).}\begin{itemize}
\item Pseudo-coherence chain. cohomology-lemma-pseudo-coherent-pullback, perfect-lemma-identify-pseudo-coherent-noetherian, cohomology-lemma-tor-pseudo-coherent (cross-chapter, named). Shows the twisted pullback is pseudo-coherent with coherent cohomology.
\item Direct image coherence. perfect-lemma-direct-image-coherent (cross-chapter, named). Reduces boundedness of the FM image to boundedness of the tensor product.
\item Tor amplitude tools. cohomology-lemma-tor-amplitude-stalk, more-algebra-lemma-complex-finite-tor-dimension-modules (cross-chapter, named). Cases (3) and (4) verify finite tor dimension stalkwise.
\end{itemize}\medskip
\ifdefined\hypertarget\hypertarget{sourcepointer-1}{}\fi
\noindent\textit{Source reference.} Original reference 00OE: focused-seeds/geometry/S9-eqv-fourier-mukai-Coh/ref-tags.txt, source lines 5--5. Navigation pointer only: the counted proof cites this statement by \textbackslash{}ref as a standing input; the deduction the counted proof has to make is carried out inside the counted span, and no part of that deduction is delegated to this pointer.\par\medskip
\ifdefined\hypertarget\hypertarget{sourcepointer-2}{}\fi
\noindent\textit{Source reference.} Original reference 09U7: focused-seeds/geometry/S9-eqv-fourier-mukai-Coh/ref-tags.txt, source lines 6--6. Navigation pointer only: the counted proof cites this statement by \textbackslash{}ref as a standing input; the deduction the counted proof has to make is carried out inside the counted span, and no part of that deduction is delegated to this pointer.\par\medskip
\ifdefined\hypertarget\hypertarget{sourcepointer-3}{}\fi
\noindent\textit{Source reference.} Original reference 09J3: focused-seeds/geometry/S9-eqv-fourier-mukai-Coh/ref-tags.txt, source lines 7--7. Navigation pointer only: the counted proof cites this statement by \textbackslash{}ref as a standing input; the deduction the counted proof has to make is carried out inside the counted span, and no part of that deduction is delegated to this pointer.\par\medskip
\ifdefined\hypertarget\hypertarget{sourcepointer-4}{}\fi
\noindent\textit{Source reference.} Original reference 09U9: focused-seeds/geometry/S9-eqv-fourier-mukai-Coh/ref-tags.txt, source lines 8--8. Navigation pointer only: the counted proof cites this statement by \textbackslash{}ref as a standing input; the deduction the counted proof has to make is carried out inside the counted span, and no part of that deduction is delegated to this pointer.\par\medskip
\ifdefined\hypertarget\hypertarget{sourcepointer-5}{}\fi
\noindent\textit{Source reference.} Original reference 08DE: focused-seeds/geometry/S9-eqv-fourier-mukai-Coh/ref-tags.txt, source lines 9--9. Navigation pointer only: the counted proof cites this statement by \textbackslash{}ref as a standing input; the deduction the counted proof has to make is carried out inside the counted span, and no part of that deduction is delegated to this pointer.\par\medskip
\ifdefined\hypertarget\hypertarget{sourcepointer-6}{}\fi
\noindent\textit{Source reference.} Original reference 066H: focused-seeds/geometry/S9-eqv-fourier-mukai-Coh/ref-tags.txt, source lines 10--10. Navigation pointer only: the counted proof cites this statement by \textbackslash{}ref as a standing input; the deduction the counted proof has to make is carried out inside the counted span, and no part of that deduction is delegated to this pointer.\par\medskip
\ifdefined\hypertarget\hypertarget{sourcepointer-7}{}\fi
\noindent\textit{Source reference.} Original reference 08E2: focused-seeds/geometry/S9-eqv-fourier-mukai-Coh/ref-tags.txt, source lines 11--11. Navigation pointer only: the counted proof cites this statement by \textbackslash{}ref as a standing input; the deduction the counted proof has to make is carried out inside the counted span, and no part of that deduction is delegated to this pointer.\par\medskip
\ifdefined\hypertarget\hypertarget{sourcepointer-8}{}\fi
\noindent\textit{Source reference.} Original reference 08E8: focused-seeds/geometry/S9-eqv-fourier-mukai-Coh/ref-tags.txt, source lines 12--12. Navigation pointer only: the counted proof cites this statement by \textbackslash{}ref as a standing input; the deduction the counted proof has to make is carried out inside the counted span, and no part of that deduction is delegated to this pointer.\par\medskip
\begin{lemma}
\label{lemma-fourier-mukai-Coh}
Let $S$ be a Noetherian scheme. Let $X$ and $Y$ be schemes of finite type
over $S$. Let $K \in D^b_{\textit{Coh}}(\mathcal{O}_{X \times_S Y})$.
The corresponding Fourier-Mukai functor $\Phi_K$ sends
$D^b_{\textit{Coh}}(\mathcal{O}_X)$ into $D^b_{\textit{Coh}}(\mathcal{O}_Y)$
if at least one of the following conditions is satisfied:
\begin{enumerate}
\item the support of $H^i(K)$ is proper over $Y$ for all $i$, and $K$
has finite tor dimension as an object of $D(\text{pr}_1^{-1}\mathcal{O}_X)$,
\item $K$ can be represented by a bounded complex $\mathcal{K}^\bullet$
of coherent $\mathcal{O}_{X \times_S Y}$-modules, flat over $X$, with support
proper over $Y$,
\item the support of $H^i(K)$ is proper over $Y$ for all $i$
and $X$ is a regular scheme,
\item $K$ is perfect, the support of $H^i(K)$ is proper over $Y$ for all $i$,
and $Y \to S$ is flat.
\end{enumerate}
Furthermore in each case the support condition is automatic
if $X \to S$ is proper.
\end{lemma}

\begin{proof}
Let $M$ be an object of $D^b_{\textit{Coh}}(\mathcal{O}_X)$.
In each case we will use Derived Categories of Schemes, Lemma
\ref{perfect-lemma-direct-image-coherent} to show that
$$
\Phi_K(M) = R\text{pr}_{2, *}(
L\text{pr}_1^*M
\otimes_{\mathcal{O}_{X \times_S Y}}^\mathbf{L}
K)
$$
is in $D^b_{\textit{Coh}}(\mathcal{O}_Y)$. The derived tensor product
$L\text{pr}_1^*M \otimes_{\mathcal{O}_{X \times_S Y}}^\mathbf{L} K$
is a pseudo-coherent object of $D(\mathcal{O}_{X \times_S Y})$
(by
Cohomology, Lemma \ref{cohomology-lemma-pseudo-coherent-pullback},
Derived Categories of Schemes, Lemma
\ref{perfect-lemma-identify-pseudo-coherent-noetherian}, and
Cohomology, Lemma \ref{cohomology-lemma-tensor-pseudo-coherent})
whence has coherent cohomology sheaves (by
Derived Categories of Schemes, Lemma
\ref{perfect-lemma-identify-pseudo-coherent-noetherian} again).
In each case the supports of the cohomology sheaves
$H^i(L\text{pr}_1^*M \otimes_{\mathcal{O}_{X \times_S Y}}^\mathbf{L} K)$
is proper over $Y$ as these supports are contained in the
union of the supports of the $H^i(K)$. Hence in each case
it suffices to prove that this tensor product is bounded below.

\medskip\noindent
Case (1). By Cohomology, Lemma \ref{cohomology-lemma-variant-derived-pullback}
we have
$$
L\text{pr}_1^*M
\otimes_{\mathcal{O}_{X \times_S Y}}^\mathbf{L}
K
\cong
\text{pr}_1^{-1}M
\otimes_{\text{pr}_1^{-1}\mathcal{O}_X}^\mathbf{L}
K
$$
with obvious notation. Hence the assumption on tor dimension
and the fact that $M$ has only a finite number of nonzero
cohomology sheaves, implies the bound we want.

\medskip\noindent
Case (2) follows because here the assumption implies that $K$ has
finite tor dimension as an object of $D(\text{pr}_1^{-1}\mathcal{O}_X)$
hence the argument in the previous paragraph applies.

\medskip\noindent
In Case (3) it is also the case that $K$ has finite tor dimension
as an object of $D(\text{pr}_1^{-1}\mathcal{O}_X)$. Namely, choose
affine opens $U = \Spec(A)$ and $V = \Spec(B)$ of $X$ and $Y$ mapping into the
affine open $W = \Spec(R)$ of $S$. Then
$K|_{U \times V}$ is given by a bounded complex of finite
$A \otimes_R B$-modules $M^\bullet$. Since $A$ is a regular ring
of finite dimension we see that each $M^i$ has finite projective dimension
as an $A$-module (Algebra, Lemma
\ref{algebra-lemma-finite-gl-dim-finite-dim-regular})
and hence finite tor dimension as an $A$-module.
Thus $M^\bullet$ has finite tor dimension as a complex of $A$-modules
(More on Algebra, Lemma
\ref{more-algebra-lemma-complex-finite-tor-dimension-modules}).
Since $X \times Y$ is quasi-compact we conclude there exist $[a, b]$
such that for every point $z \in X \times Y$ the stalk $K_z$
has tor amplitude in $[a, b]$ over $\mathcal{O}_{X, \text{pr}_1(z)}$.
This implies $K$ has bounded tor dimension as an object of
$D(\text{pr}_1^{-1}\mathcal{O}_X)$, see
Cohomology, Lemma \ref{cohomology-lemma-tor-amplitude-stalk}.
We conclude as in the previous to paragraphs.

\medskip\noindent
Case (4). With notation as above, the ring map $R \to B$ is flat.
Hence the ring map $A \to A \otimes_R B$ is flat. Hence any projective
$A \otimes_R B$-module is $A$-flat. Thus any perfect complex of
$A \otimes_R B$-modules has finite tor dimension as a complex
of $A$-modules and we conclude as before.
\end{proof}

\end{document}
